\doclearpage
\section{PIOPs for Wild-Card Pattern Matching}

In this section we develop PIOPs for SQL \emph{wild-card} pattern matching --- the predicates written with the \texttt{LIKE} operator. A \texttt{LIKE} pattern may use two kinds of wild-card: the multi-character \texttt{\%}, which matches any (possibly empty) run of characters, and the single-character \texttt{\_}, which matches exactly one character. \Cref{fig:wildcard-taxonomy} shows the taxonomy of patterns we support: three kinds --- single-character, multi-character, and mixed --- where the single- and multi-character kinds each come in a single-factor and a multi-factor mode.

\begin{figure}[H]
\centering
\resizebox{\ifsp\columnwidth\else\width\fi}{!}{%
\begin{tikzpicture}[
    level 1/.style={sibling distance=53mm, level distance=13mm},
    level 2/.style={sibling distance=25mm, level distance=13mm},
    every node/.style={draw, rounded corners, align=center, font=\normalsize, inner sep=3pt},
    edge from parent/.style={draw, semithick}
  ]
  \node {Wild-card\\pattern matching}
    child { node {Single-character}
      child { node {Single-factor} }
      child { node {Multi-factor} }
    }
    child { node {Multi-character}
      child { node {Single-factor}
      child[sibling distance=14mm] { node {Prefix} }
      child[sibling distance=14mm] { node {Suffix} }
      child[sibling distance=14mm] { node {Infix} }
      }
      child { node {Multi-factor} }
    }
    child[xshift=-20mm] { node {Mixed} };
\end{tikzpicture}%
}
\caption{Taxonomy of wild-card pattern matching}
\label{fig:wildcard-taxonomy}
\end{figure}

\subsection{Multi-Character}
\label{sec:multi-char-pattern-matching}

The most common kind of string pattern matching in natural SQL workloads is a multi-character pattern of the form $p_1 \% p_2 \% \dots \% p_t$, where each factor $p_i$ is a literal string and \% matches an arbitrary-length run of characters; the predicate holds iff the factors occur \emph{in order} and \emph{without overlap} within the value, with the leading and trailing \% (or their absence) anchoring the match to the start and/or end. A \emph{single-factor} pattern is simply the case $t=1$, so we treat it as a special case of the general protocol rather than a separate construction.


\begin{DefinePIOP}{Multi-Character Pattern Matching}{multi-char-pattern-matching}[Multi-Character Pattern Matching]
  \emergencystretch=3em\relax
  \emph{Inputs:} the pattern factors $\mathcal{P}=(p_1,\dots,p_t)$, of total length $k:=\sum_{i=1}^{t}|p_i|$, together with their anchoring modes $\text{mode}_j$; the input activators $\Oracle{\MLE{\mathsf{a^{old}}}}$ and $\Oracle{\MLE{\mathsf{char\text{-}act^{old}}}}$; the claimed result activators $\Oracle{\MLE{\mathsf{a^{new}}}}$ and $\Oracle{\MLE{\mathsf{char\text{-}act^{new}}}}$; and the fixed columns $\Oracle{\MLE{\mathsf{char}}}, \Oracle{\MLE{\mathsf{bnd}}}, \Oracle{\MLE{\mathsf{orig\text{-}ind}}}, \Oracle{\MLE{\mathsf{int\text{-}ind}}}, \Oracle{\MLE{\mathsf{ind}}}, \Oracle{\MLE{\mathsf{l}}}$.
  \vspace{2mm}
  \begin{enumerate}[nosep]
    \item (\emph{Activator validity})\label[step]{mc:validity} $\PIOPProver$ and $\PIOPVerifier$ invoke \cref{piop:data-preserving-update-check} on $(\mathsf{char\text{-}act^{new}}, \mathsf{a^{new}}, \mathsf{orig\text{-}ind}, \mathsf{ind}, \mathsf{l})$.
    \item (\emph{Length filtering})\label[step]{mc:filter} $\PIOPProver$ computes $\mathsf{a^{old'}}, \mathsf{char\text{-}act^{old'}}$, the result of dropping strings shorter than $k$, and the parties invoke \cref{piop:len-filter} on $(\mathsf{a^{old}}, \mathsf{char\text{-}act^{old}}, \mathsf{a^{old'}}, \mathsf{char\text{-}act^{old'}})$ with threshold $k$.
    \item (\emph{Containment})\label[step]{mc:containment} $\PIOPProver$ and $\PIOPVerifier$ invoke \cref{piop:zerocheck} on $\mathsf{a^{new}}\cdot(1-\mathsf{a^{old'}})$.
    \item (\emph{Sweep})\label[step]{mc:sweep} $\PIOPProver$ and $\PIOPVerifier$ invoke start sweeping patterns $p_1,\dots,p_t$ as follows:
    \begin{enumerate}[nosep]
    \item (\emph{Initialization})\label[step]{mc:init} $\PIOPProver$ and $\PIOPVerifier$ set $\mathsf{alive_h^{(0)}}:=\mathsf{a^{old'}}$ and $\mathsf{alive_c^{(0)}}:=\mathsf{char\text{-}act^{old'}}$.
        \item For $j=1,\dots,t$, $\PIOPProver$ and $\PIOPVerifier$ invoke \cref{piop:factor-placement} on $(p_j, \mathsf{alive_h^{(j-1)}},\mathsf{alive_c^{(j-1)}},\Oracle{\MLE{\mathsf{char}}}, \Oracle{\MLE{\mathsf{bnd}}}, \allowbreak\Oracle{\MLE{\mathsf{orig\text{-}ind}}},\allowbreak \Oracle{\MLE{\mathsf{int\text{-}ind}}},\allowbreak \Oracle{\MLE{\mathsf{ind}}})$ obtaining the match indicator $\mathsf{match}_j$ and placement $\mathsf{start}_j$.
    \item (\emph{Past-placement indicator})\label[step]{mc:past} $\PIOPProver$ broadcasts $\mathsf{start}_j$ to character level as $\mathsf{start}_j'$ via \cref{piop:broadcast}, and the parties invoke the Non-Negative \cref{piop:sign-check} on $\mathsf{int\text{-}ind}[c]-\mathsf{start}_{j}'[c]-|p_{j}|$, receiving the indicator we denote by $\mathsf{past}_{j}$, so that
    \begin{align*}
      \mathsf{past}_{j}[c] = \mathbf{1}\bigl[\mathsf{int\text{-}ind}[c]\ge\mathsf{start}_{j}'[c]+|p_{j}|\bigr],
    \end{align*}
    \end{enumerate} 
    \item \emph{(Updating the activators)} $\PIOPProver$ and $\PIOPVerifier$ update the following columns:
    \begin{align*}
      \mathsf{alive_h^{(j)}}[i] &:= \mathsf{alive_h^{(j-1)}}[i]\cdot\mathsf{match}_j[i]\\
      \mathsf{alive_c^{(j)}}[c] &:= \mathsf{alive_c^{(j-1)}}[c]\cdot\mathsf{alive_h^{(j)}}[\mathsf{orig\text{-}ind}[c]]\cdot\mathsf{past}_{j}[c],
    \end{align*}
      Then $\PIOPProver$ broadcasts $\mathsf{alive_h^{(j)}}$ to character level $\mathsf{alive_h^{(j)}}'$, and the parties invoke a \cref{piop:boolean-check} on $\mathsf{alive_c^{(j)}}$ and a \cref{piop:zerocheck} on $\mathsf{alive_c^{(j)}}-\mathsf{alive_c^{(j-1)}}\cdot\mathsf{alive_h^{(j)}}'\cdot\mathsf{past}_{j}$.
    \item (\emph{Verdict})\label[step]{mc:verdict} $\PIOPProver$ and $\PIOPVerifier$ invoke \cref{piop:zerocheck} on $\mathsf{a^{new}}\cdot(1-\mathsf{alive_h^{(t)}})$ and on $\mathsf{alive_h^{(t)}}\cdot(1-\mathsf{a^{new}})$, certifying $\mathsf{a^{new}}=\mathsf{alive_h^{(t)}}$.
  \end{enumerate}
\end{DefinePIOP}

the pattern $p=(\mathsf{str},\mathsf{mode})$ where $\mathsf{str}$ of length $\ell$ and a mode $\in\{\text{prefix},\text{suffix},\text{infix}\}$; the input activators $\Oracle{\MLE{\mathsf{a}}}$  and $\Oracle{\MLE{\mathsf{char\text{-}act}}}$  delimiting the active region; and the fixed columns $\Oracle{\MLE{\mathsf{char}}}, \Oracle{\MLE{\mathsf{bnd}}}, \Oracle{\MLE{\mathsf{orig\text{-}ind}}}, \Oracle{\MLE{\mathsf{int\text{-}ind}}}, \Oracle{\MLE{\mathsf{ind}}}$.

\begin{DefinePIOP}{Factor Placement}{factor-placement}[Factor Placement]
  \emph{Inputs:} $p=(\mathsf{str},\mathsf{mode}), \Oracle{\MLE{\mathsf{a}}},\Oracle{\MLE{\mathsf{char\text{-}act}}},\Oracle{\MLE{\mathsf{char}}}, \Oracle{\MLE{\mathsf{bnd}}}, \Oracle{\MLE{\mathsf{orig\text{-}ind}}}, \Oracle{\MLE{\mathsf{int\text{-}ind}}}, \Oracle{\MLE{\mathsf{ind}}}$.
  \vspace{2mm}
  \begin{enumerate}[nosep]
    \item (\textit{Rotated character columns}) $\PIOPProver$ computes $\mathsf{char}^{(1)},\ldots,\mathsf{char}^{(\ell-1)}$ and sends them to $\PIOPVerifier$, where each $\mathsf{char}^{(\delta)}$ is a rotation of $\mathsf{char}^{(\delta-1)}$ by $-1$ (so $\mathsf{char}^{(\delta)}[p] = \mathsf{char}[p+\delta]$) and $\mathsf{char}^{(0)} := \mathsf{char}$.  Then for $\delta=1,\dots,\ell-1$, the parties invoke \cref{piop:rotcheck} on each pair $(\mathsf{char}^{(\delta-1)}, \mathsf{char}^{(\delta)}, -1)$.
    \item (\textit{Window fingerprinting}) $\PIOPVerifier$ samples $r_0,\ldots,r_{\ell-1} \samples \F$; the parties form $\mathsf{wf} := \sum_{\delta=0}^{\ell-1} r_\delta\,\mathsf{char}^{(\delta)}$, $\mathsf{pf} := \sum_{\delta=0}^{\ell-1} r_\delta\,\mathsf{str}[\delta]$, and $\mathsf{diff} := \mathsf{wf}-\mathsf{pf}$.
    \item (\textit{Attention mask}) $\PIOPProver$ and $\PIOPVerifier$ fix the attention mask $\mathsf{att\_mask}$ according to the mode:
    \begin{enumerate}[nosep]
      \item \emph{prefix}: $\mathsf{att\_mask} := \mathsf{char\text{-}act}\cdot\mathsf{bnd}$;
      \item \emph{suffix}: $\PIOPProver$ computes $\mathsf{att\_mask} := \rho_{-\ell}(\mathsf{char\text{-}act}\cdot\mathsf{bnd})$ and sends it to $\PIOPVerifier$. The parties then invoke \cref{piop:rotcheck} on $(\mathsf{char\text{-}act}\cdot\mathsf{bnd},\,\mathsf{att\_mask},\,-\ell)$;
      \item \emph{infix}: $\PIOPProver$ computes the rotated boundary markers $\mathsf{bnd}^{(1)},\ldots,\mathsf{bnd}^{(\ell-1)}$ and sends them to $\PIOPVerifier$, where $\mathsf{bnd}^{(0)}:=\mathsf{bnd}$ and each $\mathsf{bnd}^{(j)}$ is a rotation of $\mathsf{bnd}^{(j-1)}$ by $-1$ (so $\mathsf{bnd}^{(j)}[i] = \mathsf{bnd}[i+j]$); the parties invoke \cref{piop:rotcheck} for each $(\mathsf{bnd}^{(j-1)},\mathsf{bnd}^{(j)},-1)$ and set $\mathsf{att\_mask} := \mathsf{char\text{-}act}\cdot\bigl(1-\textstyle\sum_{j=1}^{\ell-1}\mathsf{bnd}^{(j)}\bigr)$.
    \end{enumerate}
    \item (\textit{Occurrence column})\label[step]{fp:occurs} $\PIOPProver$ computes $\mathsf{occurs}$ and sends it to $\PIOPVerifier$, which is $1$ at the start of every occurrence of $\mathsf{str}$ among the candidate windows, and the parties establish its validity by
    \begin{enumerate}[nosep]
      \item a \cref{piop:boolean-check} on $\mathsf{occurs}$;
      \item a \cref{piop:zerocheck} on $\mathsf{occurs}\cdot(1-\mathsf{att\_mask})$, so every mark is a candidate window;
      \item a \cref{piop:zerocheck} on $\mathsf{occurs}\cdot\mathsf{diff}$, so every marked window equals $\mathsf{str}$;
      \item a \cref{piop:nozero-check} on the column with data $\mathsf{diff}$ and activator $\mathsf{att\_mask}\cdot(1-\mathsf{occurs})$, so every unmarked candidate differs from $\mathsf{str}$,
    \end{enumerate}
    forcing $\mathsf{occurs}$ to mark exactly the candidate windows equal to $\mathsf{str}$.
    \item (\textit{Leftmost placement})\label[step]{fp:placement} Writing $O_i:=\{\,\mathsf{int\text{-}ind}[c] : \mathsf{orig\text{-}ind}[c]=i,\ \mathsf{occurs}[c]=1\,\}$ for the occurrence positions in string $i$, $\PIOPProver$ builds and sends
    \begin{itemize}[nosep]
      \item the string-level indicator $\mathsf{match}[i]:=\mathbf{1}[\,O_i\neq\emptyset\,]$ (the string carries an occurrence);
      \item the word-level placement $\mathsf{start}[i]:=\min O_i$ when $\mathsf{match}[i]=1$ (the leftmost occurrence), and $\mathsf{start}[i]:=0$ otherwise;
      \item the character-level marker $\mathsf{mark}[c]:=\mathbf{1}[\,\mathsf{occurs}[c]=1\ \text{and}\ \mathsf{int\text{-}ind}[c]=\mathsf{start}[\mathsf{orig\text{-}ind}[c]]\,]$ (the character at that leftmost occurrence).
    \end{itemize}
    The parties then check
    \begin{enumerate}[nosep]
      \item a \cref{piop:boolean-check} on $\mathsf{match}$ and on $\mathsf{mark}$;
      \item (\emph{no false negatives}) $\PIOPProver$ broadcasts $\mathsf{match}$ to character level ($\mathsf{match}'$) via \cref{piop:broadcast}, and the parties invoke a \cref{piop:zerocheck} on $\mathsf{occurs}\cdot(1-\mathsf{match}')$, so every occurrence lies in a matched string;
      \item (\emph{no false positives}) a \cref{piop:zerocheck} on $\mathsf{mark}\cdot(1-\mathsf{occurs})$ (each mark is an occurrence), a \cref{piop:nodup-check} on $\mathsf{orig\text{-}ind}$ with activator $\mathsf{mark}$ (at most one mark per string), and---with a count $n_{\mathsf{m}}$ from $\PIOPProver$---two \cref{piop:sumcheck}s for $\sum_c\mathsf{mark}[c]=n_{\mathsf{m}}$ and $\sum_i\mathsf{match}[i]=n_{\mathsf{m}}$, so exactly one occurrence is marked in each matched string;
      \item (\emph{placement}) $\PIOPVerifier$ samples $r\samples\F$ and the parties invoke \cref{piop:lookup-check} on $(\mathsf{match},\ \mathsf{ind}+r\,\mathsf{start}) \lookup (\mathsf{mark},\ \mathsf{orig\text{-}ind}+r\,\mathsf{int\text{-}ind})$, i.e.\ with $\mathsf{match}$ activating the string-level column $\mathsf{ind}+r\,\mathsf{start}$ and $\mathsf{mark}$ activating the character-level column $\mathsf{orig\text{-}ind}+r\,\mathsf{int\text{-}ind}$. This ties each matched string's $\mathsf{start}$ to the $\mathsf{int\text{-}ind}$ of its mark (both activated sides have $n_{\mathsf{m}}$ elements, so the lookup forces equality);
      \item (\emph{leftmost}) $\PIOPProver$ broadcasts $\mathsf{start}$ to character level ($\mathsf{start}'$); a \cref{piop:sign-check} forms $\mathbf{1}[\mathsf{int\text{-}ind}<\mathsf{start}']$, and the parties invoke a \cref{piop:zerocheck} on $\mathsf{occurs}\cdot\mathbf{1}[\mathsf{int\text{-}ind}<\mathsf{start}']$ with activator $\mathsf{match}'$, so no occurrence precedes the marked one.
    \end{enumerate}
  \end{enumerate}
\end{DefinePIOP}

\ifsp\begin{figure*}[p]\else\begin{figure}[p]\fi
\centering
\resizebox{\textwidth}{!}{%
\begin{tikzpicture}[
    cell/.style={draw, minimum size=5.6mm, inner sep=0pt, font=\small},
    scell/.style={draw, minimum size=5.6mm, inner sep=0pt, font=\small},
    rowlab/.style={font=\small\itshape, anchor=east},
    hd/.style={font=\small\bfseries}]
  % Sweep example: pattern ab%bd%ef, factors ab, bd, ef; words:
  %   w0 ghiabkbd (drops R3), w1 abdabf (drops R2), w2 abbdef (INACTIVE), w3 abxbdgefk (matches).
  \def\cw{0.58}
  % ---- y coordinates ----
  \def\ychar{0}\def\yorig{-0.6}\def\yint{-1.2}\def\ybnd{-1.8}\def\yacold{-2.4}\def\yac{-3.0}
  \def\yocA{-5.1}\def\yaA{-5.7}
  \def\yocB{-8.3}\def\yaB{-8.9}
  \def\yocC{-11.5}\def\yaC{-12.1}
  % ---- word-name headers over the character groups ----
  \foreach \cx/\nm in {2.03/ghiabkbd, 5.51/abbd, 8.41/abdabf, 11.89/abbdef, 16.24/abxbdgefk}{
    \node[font=\footnotesize\ttfamily] at (\cx,0.92){\nm};}
  % ---- base rows: char / orig-ind / int-ind / bnd ----
  \foreach \p/\ch/\o/\ii/\bb in {%
      0/g/0/0/1,1/h/0/1/0,2/i/0/2/0,3/a/0/3/0,4/b/0/4/0,5/k/0/5/0,6/b/0/6/0,7/d/0/7/0,
      8/a/1/0/1,9/b/1/1/0,10/b/1/2/0,11/d/1/3/0,
      12/a/2/0/1,13/b/2/1/0,14/d/2/2/0,15/a/2/3/0,16/b/2/4/0,17/f/2/5/0,
      18/a/3/0/1,19/b/3/1/0,20/b/3/2/0,21/d/3/3/0,22/e/3/4/0,23/f/3/5/0,
      24/a/4/0/1,25/b/4/1/0,26/x/4/2/0,27/b/4/3/0,28/d/4/4/0,29/g/4/5/0,30/e/4/6/0,31/f/4/7/0,32/k/4/8/0}{
    \pgfmathsetmacro{\px}{\p*\cw}
    \edef\gc{\ifcase\o blue!15\or purple!16\or green!18\or black!12\or orange!22\fi}
    \edef\bc{\ifnum\bb=1 yellow!35\else gray!8\fi}
    \node[font=\scriptsize,gray] at (\px,0.52){\p};
    \node[cell,fill=\gc]    at (\px,\ychar){\texttt{\ch}};
    \node[cell,fill=\gc]    at (\px,\yorig){\o};
    \node[cell,fill=gray!6] at (\px,\yint){\ii};
    \node[cell,fill=\bc]    at (\px,\ybnd){\bb};}
  % ---- char-act^{old} (pre-filter): active on w0,abbd,w2,w4; off only on inactive abbdef (w3) ----
  \foreach \p in {0,...,32}{\pgfmathsetmacro{\px}{\p*\cw}\node[cell,fill=gray!14] at (\px,\yacold){0};}
  \foreach \p in {0,...,17}{\pgfmathsetmacro{\px}{\p*\cw}\node[cell,fill=green!30] at (\px,\yacold){1};}
  \foreach \p in {24,...,32}{\pgfmathsetmacro{\px}{\p*\cw}\node[cell,fill=green!30] at (\px,\yacold){1};}
  % ---- char-act^{old'} = alive_c^{(0)} (post-filter): active on w0,w2,w4; off on filtered abbd (w1) and inactive abbdef (w3) ----
  \foreach \p in {0,...,32}{\pgfmathsetmacro{\px}{\p*\cw}\node[cell,fill=gray!14] at (\px,\yac){0};}
  \foreach \p in {0,...,7}{\pgfmathsetmacro{\px}{\p*\cw}\node[cell,fill=green!30] at (\px,\yac){1};}
  \foreach \p in {12,...,17}{\pgfmathsetmacro{\px}{\p*\cw}\node[cell,fill=green!30] at (\px,\yac){1};}
  \foreach \p in {24,...,32}{\pgfmathsetmacro{\px}{\p*\cw}\node[cell,fill=green!30] at (\px,\yac){1};}
  % ==================== Round 1 : p1 = ab ====================
  \draw[gray!55] (-2.6,-3.8) -- (23.0,-3.8);
  \node[hd,anchor=west] at (-2.6,-4.3){Round 1: place $p_1=\texttt{ab}$};
  \foreach \p in {0,...,32}{\pgfmathsetmacro{\px}{\p*\cw}\node[cell,fill=gray!8] at (\px,\yocA){0};}
  \foreach \p in {3,12,15,24}{\pgfmathsetmacro{\px}{\p*\cw}\node[cell,fill=orange!45] at (\px,\yocA){1};}
  \foreach \p in {0,...,32}{\pgfmathsetmacro{\px}{\p*\cw}\node[cell,fill=gray!14] at (\px,\yaA){0};}
  \foreach \p in {5,6,7,14,15,16,17,26,27,28,29,30,31,32}{\pgfmathsetmacro{\px}{\p*\cw}\node[cell,fill=green!30] at (\px,\yaA){1};}
  % ==================== Round 2 : p2 = bd ====================
  \draw[gray!55] (-2.6,-7.0) -- (23.0,-7.0);
  \node[hd,anchor=west] at (-2.6,-7.5){Round 2: place $p_2=\texttt{bd}$};
  \foreach \p in {0,...,32}{\pgfmathsetmacro{\px}{\p*\cw}\node[cell,fill=gray!8] at (\px,\yocB){0};}
  \foreach \p in {6,27}{\pgfmathsetmacro{\px}{\p*\cw}\node[cell,fill=orange!45] at (\px,\yocB){1};}
  \foreach \p in {0,...,32}{\pgfmathsetmacro{\px}{\p*\cw}\node[cell,fill=gray!14] at (\px,\yaB){0};}
  \foreach \p in {29,30,31,32}{\pgfmathsetmacro{\px}{\p*\cw}\node[cell,fill=green!30] at (\px,\yaB){1};}
  % ==================== Round 3 : p3 = ef ====================
  \draw[gray!55] (-2.6,-10.2) -- (23.0,-10.2);
  \node[hd,anchor=west] at (-2.6,-10.7){Round 3: place $p_3=\texttt{ef}$};
  \foreach \p in {0,...,32}{\pgfmathsetmacro{\px}{\p*\cw}\node[cell,fill=gray!8] at (\px,\yocC){0};}
  \foreach \p in {30}{\pgfmathsetmacro{\px}{\p*\cw}\node[cell,fill=orange!45] at (\px,\yocC){1};}
  \foreach \p in {0,...,32}{\pgfmathsetmacro{\px}{\p*\cw}\node[cell,fill=gray!14] at (\px,\yaC){0};}
  \foreach \p in {32}{\pgfmathsetmacro{\px}{\p*\cw}\node[cell,fill=green!30] at (\px,\yaC){1};}
  % ---- placement brackets on the char row of the surviving word w4 (ab|bd|ef) ----
  \draw[black,thick,rounded corners=1pt] (24*\cw-0.30,0.30) rectangle (25*\cw+0.30,-0.30);
  \draw[blue!70,thick,rounded corners=1pt] (27*\cw-0.30,0.30) rectangle (28*\cw+0.30,-0.30);
  \draw[red!70,thick,rounded corners=1pt] (30*\cw-0.30,0.30) rectangle (31*\cw+0.30,-0.30);
  % ---- character-panel row labels ----
  \node[rowlab] at (-0.45,\ychar){$\mathsf{char}$};
  \node[rowlab] at (-0.45,\yorig){$\mathsf{orig\text{-}ind}$};
  \node[rowlab] at (-0.45,\yint){$\mathsf{int\text{-}ind}$};
  \node[rowlab] at (-0.45,\ybnd){$\mathsf{bnd}$};
  \node[rowlab] at (-0.45,\yacold){$\mathsf{char\text{-}act^{old}}$};
  \node[rowlab] at (-0.45,\yac){$\mathsf{char\text{-}act^{old'}}$};
  \node[font=\scriptsize\itshape,anchor=east] at (-0.45,-3.42){$=\mathsf{alive_c^{(0)}}$};
  \node[rowlab] at (-0.45,\yocA){$\mathsf{occurs}$};
  \node[rowlab] at (-0.45,\yaA){$\mathsf{alive_c^{(1)}}$};
  \node[rowlab] at (-0.45,\yocB){$\mathsf{occurs}$};
  \node[rowlab] at (-0.45,\yaB){$\mathsf{alive_c^{(2)}}$};
  \node[rowlab] at (-0.45,\yocC){$\mathsf{occurs}$};
  \node[rowlab] at (-0.45,\yaC){$\mathsf{alive_c^{(3)}}$};
  %% ============ string-level columns, one 5-word block to the right of each round band ============
  \def\lblx{20.9}
  \def\colA{21.55}\def\colB{22.13}\def\colC{22.71}
  \newcommand{\SB}[4]{\node[scell,fill=#3] at (#1,#2){#4};}
  % vertical divider separating the character-level panel from the string-level blocks
  \draw[gray!70,thick] (19.15,1.2) -- (19.15,-13.4);
  % ===== input block (beside the base character rows); shows length filtering a^old -> a^old' =====
  \foreach \yy/\nm in {0/ghiabkbd,-0.6/abbd,-1.2/abdabf,-1.8/abbdef,-2.4/abxbdgefk}{
    \node[font=\scriptsize\ttfamily,anchor=east] at (\lblx,\yy){\nm};}
  % column headers, once at the top (input meaning / round-$j$ meaning); columns align across all blocks
  \foreach \c/\lab in {\colA/{$\mathsf{l}\,/\,\mathsf{match}_j$},\colB/{$\mathsf{a^{old}}\,/\,\mathsf{start}_j$},\colC/{$\mathsf{a^{old'}}\,/\,\mathsf{alive_h^{(j)}}$}}{
    \node[font=\scriptsize,rotate=90,anchor=west] at (\c,0.4){\lab};}
  \SB{\colA}{0}{gray!8}{8}\SB{\colB}{0}{green!30}{1}\SB{\colC}{0}{green!30}{1}
  \SB{\colA}{-0.6}{gray!8}{4}\SB{\colB}{-0.6}{green!30}{1}\SB{\colC}{-0.6}{red!30}{0}
  \SB{\colA}{-1.2}{gray!8}{6}\SB{\colB}{-1.2}{green!30}{1}\SB{\colC}{-1.2}{green!30}{1}
  \SB{\colA}{-1.8}{gray!8}{6}\SB{\colB}{-1.8}{gray!15}{0}\SB{\colC}{-1.8}{gray!15}{0}
  \SB{\colA}{-2.4}{gray!8}{9}\SB{\colB}{-2.4}{green!30}{1}\SB{\colC}{-2.4}{green!30}{1}
  % ===== Round 1 block =====
  \foreach \yy/\nm in {-4.2/ghiabkbd,-4.8/abbd,-5.4/abdabf,-6.0/abbdef,-6.6/abxbdgefk}{
    \node[font=\scriptsize\ttfamily,anchor=east] at (\lblx,\yy){\nm};}
  \SB{\colA}{-4.2}{green!30}{1}\SB{\colB}{-4.2}{gray!8}{3}\SB{\colC}{-4.2}{green!30}{1}
  \SB{\colA}{-4.8}{gray!14}{--}\SB{\colB}{-4.8}{gray!14}{--}\SB{\colC}{-4.8}{gray!15}{0}
  \SB{\colA}{-5.4}{green!30}{1}\SB{\colB}{-5.4}{gray!8}{0}\SB{\colC}{-5.4}{green!30}{1}
  \SB{\colA}{-6.0}{gray!14}{--}\SB{\colB}{-6.0}{gray!14}{--}\SB{\colC}{-6.0}{gray!15}{0}
  \SB{\colA}{-6.6}{green!30}{1}\SB{\colB}{-6.6}{gray!8}{0}\SB{\colC}{-6.6}{green!30}{1}
  % ===== Round 2 block =====
  \foreach \yy/\nm in {-7.4/ghiabkbd,-8.0/abbd,-8.6/abdabf,-9.2/abbdef,-9.8/abxbdgefk}{
    \node[font=\scriptsize\ttfamily,anchor=east] at (\lblx,\yy){\nm};}
  \SB{\colA}{-7.4}{green!30}{1}\SB{\colB}{-7.4}{gray!8}{6}\SB{\colC}{-7.4}{green!30}{1}
  \SB{\colA}{-8.0}{gray!14}{--}\SB{\colB}{-8.0}{gray!14}{--}\SB{\colC}{-8.0}{gray!15}{0}
  \SB{\colA}{-8.6}{red!30}{0}\SB{\colB}{-8.6}{gray!14}{--}\SB{\colC}{-8.6}{gray!15}{0}
  \SB{\colA}{-9.2}{gray!14}{--}\SB{\colB}{-9.2}{gray!14}{--}\SB{\colC}{-9.2}{gray!15}{0}
  \SB{\colA}{-9.8}{green!30}{1}\SB{\colB}{-9.8}{gray!8}{3}\SB{\colC}{-9.8}{green!30}{1}
  % ===== Round 3 block; alive_h^{(3)} is the final output a^new =====
  \foreach \yy/\nm in {-10.6/ghiabkbd,-11.2/abbd,-11.8/abdabf,-12.4/abbdef,-13.0/abxbdgefk}{
    \node[font=\scriptsize\ttfamily,anchor=east] at (\lblx,\yy){\nm};}
  \SB{\colA}{-10.6}{red!30}{0}\SB{\colB}{-10.6}{gray!14}{--}\SB{\colC}{-10.6}{gray!15}{0}
  \SB{\colA}{-11.2}{gray!14}{--}\SB{\colB}{-11.2}{gray!14}{--}\SB{\colC}{-11.2}{gray!15}{0}
  \SB{\colA}{-11.8}{gray!14}{--}\SB{\colB}{-11.8}{gray!14}{--}\SB{\colC}{-11.8}{gray!15}{0}
  \SB{\colA}{-12.4}{gray!14}{--}\SB{\colB}{-12.4}{gray!14}{--}\SB{\colC}{-12.4}{gray!15}{0}
  \SB{\colA}{-13.0}{green!30}{1}\SB{\colB}{-13.0}{gray!8}{6}\SB{\colC}{-13.0}{green!45}{1}
  % closing line under round 3 (alive_h^{(3)} is the final output a^new; see caption)
  \draw[gray!55] (-2.6,-13.4) -- (23.0,-13.4);
\end{tikzpicture}}
\caption{Worked example of the \textsc{Multi-Character Pattern Matching} sweep (\cref{piop:multi-char-pattern-matching}) on \texttt{LIKE 'ab\%bd\%ef'} (factors $\mathcal{P}=(\texttt{ab},\texttt{bd},\texttt{ef})$). The character-level rows span all $33$ positions; to the right of each round's character rows sits that round's string-level block, one row per word, so every sweep round gathers its character- and string-level polynomials together. Word \texttt{abbdef} is inactive from the start, and \texttt{abbd} is removed by length filtering ($\mathsf{a^{old}}=1$ but $\mathsf{a^{old'}}=0$, as its $4$ characters fall short of $k=6$), so neither enters the sweep; the $\mathsf{char\text{-}act^{old}}$ and $\mathsf{char\text{-}act^{old'}}$ rows show the same filtering at character level, the latter being $\mathsf{alive_c^{(0)}}$ that seeds the sweep. Round~1 places \texttt{ab} in every surviving word (all stay alive); Round~2 drops \texttt{abdabf}, which has no \texttt{bd} past its \texttt{ab}; Round~3 drops \texttt{ghiabkbd}, whose suffix is exhausted before an \texttt{ef}. Only \texttt{abxbdgefk} survives, so $\mathsf{a^{new}}=\mathsf{alive_h^{(3)}}$ keeps it alone; its in-order embedding \texttt{ab}$\mid$\texttt{bd}$\mid$\texttt{ef} is boxed on the character row. The $\mathsf{alive_c^{(j)}}$ rows show the working suffix of each word shrinking as factors are consumed.}
\label{fig:sweep-example}
\ifsp\end{figure*}\else\end{figure}\fi

